\documentclass[../../../include/open-logic-section]{subfiles}

\begin{document}

\olfileid{sth}{z}{nat}
\olsection{આપણી પ્રાકૃતિક સંખ્યાઓની પસંદગી}

%Let us take it for granted that the informal principle \stagesinf{}
%can be justified. It is also, though, worth commenting on the
%particular axiom we extracted from that informal principle. 

\olref[infinity-again]{defnomega}માં આપણે “પ્રાકૃતિક સંખ્યાઓ”
શબ્દપ્રયોગની સ્પષ્ટ વ્યાખ્યા આપી. આ ઠરાવ કેવી રીતે સમજવો?
તે કોઈ તત્ત્વમીમાંસક દાવો નથી, પણ અમુક ગણોને પ્રાકૃતિક સંખ્યાઓ
\emph{તરીકે લેવાનો} નિર્ણય જ છે. આવું માનવાનાં કારણોનો થોડો
ઉલ્લેખ \olref[sfr][rel][ref]{sec}, \olref[sfr][arith][ref]{sec}
અને \olref[sfr][infinite][dedekindsproof]{sec}માં કર્યો હતો.
પણ આ કારણોને વધુ સ્પષ્ટ રીતે રજૂ કરી શકાય.

આપણી અનંતતાની સ્વયંસિદ્ધિ \citet{VonNeumann1925}ને અનુસરે છે.
પણ તેના બદલે સ્વીકારી શક્યા હોત એવી બીજી સ્વયંસિદ્ધિ આ છે:

\begin{defish}
\emph{ઝર્મેલોની \citeyear{Zermelo1908Untersuchungen}ની અનંતતાની
સ્વયંસિદ્ધિ.} એવો ગણ $A$ છે કે $\emptyset \in A$ અને
$(\forall x \in A)\{x\} \in A$.
\end{defish}

આપણી (વૉન નોયમનથી પ્રેરિત) અનંતતાની સ્વયંસિદ્ધિને બદલે
ઝર્મેલોની સ્વયંસિદ્ધિ લીધી હોત, તો પણ એટલી જ સારી રીતે
ડેડેકિન્ડ-અનંત ગણ અને તેથી ડેડેકિન્ડ બીજગણિત મળ્યાં હોત.
ઝર્મેલોના અભિગમમાં આપણા બીજગણિતનો વિશિષ્ટ ઘટક ફરી
$\emptyset$ (સંખ્યા $0$નું આપણું અવેજ) જ હોત; પરંતુ એક-એક
વિધેય $x \mapsto x \cup \{x\}$ને બદલે
$x \mapsto \{x\}$ નકશા દ્વારા અપાત. આનો સૌથી સરળ
ફલિતાર્થ એ છે કે ઝર્મેલો $2$ને $\{\{\emptyset\}\}$ ગણે છે,
જ્યારે આપણે (વૉન નોયમન સાથે) $2$ને
$\{\emptyset, \{\emptyset\}\}$ ગણીએ છીએ.

અનંતતાની એક સ્વયંસિદ્ધિ બીજી કરતાં શા માટે પસંદ કરવી?
મુખ્ય વ્યવહારુ કારણ એ છે કે વૉન નોયમનનો અભિગમ મોટા પાયે
વિસ્તારીને અનંતાતીત સંખ્યાઓને સારી રીતે સંભાળી શકે છે.
\olref[ordinals][]{chap}થી આગળ આનો અભ્યાસ કરીશું.
જોકે માત્ર \emph{અંકગણિત કરવાની} દૃષ્ટિએ બંને અભિગમો
એટલી જ સારી રીતે કામ કરે. તેથી કોઈ કહે કે પ્રાકૃતિક
સંખ્યાઓ ગણો \emph{છે}, તો સ્વાભાવિક પ્રશ્ન છે:
\emph{તે કયા ગણો છે?}

બરાબર આ પ્રશ્નને \citet{Benacerraf1965}એ પ્રસિદ્ધ કર્યો.
પણ તે આપણે અગાઉ ઘણી વાર જોયેલી ઘટનાનું સૌથી પ્રસિદ્ધ
ઉદાહરણ જ છે તે પર ભાર મૂકવો જોઈએ. મૂળ મુદ્દો આ છે.
ગણસિદ્ધાંત આપણને અનેક “સહજ” પ્રકારની વસ્તુઓનું
\emph{અનુકરણ} કરવાની રીત આપે છે: વાસ્તવિક સંખ્યાઓ,
સંમેય સંખ્યાઓ, પૂર્ણાંકો અને પ્રાકૃતિક સંખ્યાઓ તો ખરાં
જ; પણ ક્રમયુક્ત જોડ, વિધેયો અને સંબંધો પણ.
જોકે ગણસિદ્ધાંત અનુકરણની \emph{અનન્ય} પસંદગી ક્યારેય
આપતો નથી. આપણા હેતુ માટે સીધી રીતે એટલું જ સારું
કામ કર્યું હોત એવા વિકલ્પો \emph{હંમેશાં} હોય છે.

\end{document}
